
\section{Cytoscape}
\label{research:Cytoscape}
\textbf{Source:} \cite{Cytoscapejs}
Once I had settled on the idea of displaying my note structure using a mindmap-style graph, I began searching for technologies that would allow me to create visually appealing and functional diagrams.

I initially explored tools such as:
\begin{itemize}
    \item GoJS
    \item SciChart
    \item ReactFlow
\end{itemize}
 While each showed promise, they quickly revealed some critical limitations — lack of flexibility, poor integration with frontend libraries, and a general inability to produce the kind of polished, aesthetic visuals I was aiming for. This prompted me to shift my focus away from browser-based libraries and instead look into more complete desktop applications that already offer the features I was after.

That’s when I came across Cytoscape.js, a powerful graphing tool originally developed for visualising biological data relationships. Despite its scientific roots, Cytoscape.js turned out to be a strong match for my project. I was especially pleased to find that it came with a JavaScript library that could be integrated directly into my app, allowing me to create high-quality, interactive graphs with minimal setup.

Another key advantage was the active community built around the tool. This community has developed a variety of plugins — including right-click context menus and, more importantly, multiple layout ‘engines’. These layout options are essential to my goal of generating diagrams that are both functional and visually balanced, and I plan to explore them further as I refine the graphical side of my project.

I spent some time reading through the documentation to learn more about the program, and then created a prototype proof of concept which I am very pleased with
